\documentclass[../../../hsm_tok_q.tex]{subfiles}

\begin{document}
    \setcounter{figure}{0}
    \textsf{図\ref{fig/hsm_tok_2025_1st_04_01_q}}で，%
    点Oは線分ABを直径とする半円の中心である．

    点Pは，線分OA上にある点で，点O，点Aのいずれにも一致しない．

    点Qは，$\myarca{\mathrm{AB}}$上にある点で，点A，点Bのいずれにも一致しない．

    点Rは，$\myarca{\mathrm{BQ}}$上にある点で，点B，点Qのいずれにも一致しない．

    点Aと点Q，点Aと点R，点Bと点Q，点Pと点Rをそれぞれ結ぶ．
    
    次の各問に答えよ．
    \\
    \begin{minipage}{\linewidth}
        \vspace{.5\baselineskip}
        {\centering
            \setlength{\abovecaptionskip}{3pt}
            \includegraphics%
                {\subfix{fig_hsm_tok_2025_1st_04/fig_hsm_tok_2025_1st_04_01_q.pdf}}%
                \captionof{figure}{} \label{fig/hsm_tok_2025_1st_04_01_q}%
        }%
    \end{minipage}
    \vspace{.1\baselineskip}

    \begin{enumerate}
        \item \textsf{図\ref{fig/hsm_tok_2025_1st_04_01_q}}において，\\
            $\mathrm{AQ}\jequal\mathrm{BQ}$，$\angle\mathrm{QAR}\jequal 20^\circ$，%
            $\angle\mathrm{ARP}\jequal a^\circ$とするとき，\\
            $\angle\mathrm{BPR}$の大きさを表す式を，%
            次の\textsf{ア}～\textsf{エ}のうちから選び，記号で答えよ．%

            {\renewcommand{\arraystretch}{1.3}%
                \begin{tabularx}{\dimexpr\linewidth-3.5mm}{@{}LLLL@{}}
                    {\textsf ア}\; $(a+20)$度 & {\textsf イ}\; $(a+25)$度 \\
                    {\textsf ウ}\; $(155-a)$度 & {\textsf エ}\; $(160-a)$度
                \end{tabularx}
            }
        \newpage
        \item \textsf{図\ref{fig/hsm_tok_2025_1st_04_02_q}}は，%
            \textsf{図\ref{fig/hsm_tok_2025_1st_04_01_q}}において，%
            $\mathrm{AP}\jequal\mathrm{AQ}$，$\myarca{\mathrm{BR}}\jequal\myarca{\mathrm{QR}}$のとき，%
            点Qと点Rを結んだ場合を表している．

            次の\ref{hsm_tok_2025_1st_04_q:1}，\ref{hsm_tok_2025_1st_04_q:2}に答えよ．
            
            \begin{myenub}
                \item $\triangle\mathrm{APR}\jequiv\triangle\mathrm{AQR}$であることを証明せよ．%
                    \label{hsm_tok_2025_1st_04_q:1}

                \item 次の\:\myboxa{　}\:の中の「\textsf{う}」「\textsf{え}」%
                    に当てはまる数字をそれぞれ答えよ．%
                    \label{hsm_tok_2025_1st_04_q:2}

                    \textsf{図\ref{fig/hsm_tok_2025_1st_04_02_q}}において，%
                    線分ARと線分BQとの交点をS，点Oと点Rを結び，%
                    線分BQと線分ORとの交点をTとした場合を考える．

                    $\mathrm{AP}\jequal 2\,\mathrm{OP}$ のとき，$\triangle\mathrm{RST}$の面積は，\\
                    四角形AORQの面積の $\myfracc{\myboxf{う}}{\myboxf{え}}$ 倍である．
            \end{myenub}                      
            \begin{minipage}{\linewidth}
                \vspace{.5\baselineskip}
                {\centering
                    \setlength{\abovecaptionskip}{3pt}
                    \includegraphics%
                        {\subfix{fig_hsm_tok_2025_1st_04/fig_hsm_tok_2025_1st_04_02_q.pdf}}%
                        \captionof{figure}{} \label{fig/hsm_tok_2025_1st_04_02_q}%
                }%
            \end{minipage}            
        \end{enumerate}

\end{document}